% Options for packages loaded elsewhere
\PassOptionsToPackage{unicode}{hyperref}
\PassOptionsToPackage{hyphens}{url}
\documentclass[
  10pt,
  titlepage]{article}
\usepackage{xcolor}
\usepackage[a4paper, includeheadfoot, margin=2.54cm]{geometry}
\usepackage{amsmath,amssymb}
\setcounter{secnumdepth}{-\maxdimen} % remove section numbering
\usepackage{iftex}
\ifPDFTeX
  \usepackage[T1]{fontenc}
  \usepackage[utf8]{inputenc}
  \usepackage{textcomp} % provide euro and other symbols
\else % if luatex or xetex
  \usepackage{unicode-math} % this also loads fontspec
  \defaultfontfeatures{Scale=MatchLowercase}
  \defaultfontfeatures[\rmfamily]{Ligatures=TeX,Scale=1}
\fi
\usepackage{lmodern}
\ifPDFTeX\else
  % xetex/luatex font selection
\fi
% Use upquote if available, for straight quotes in verbatim environments
\IfFileExists{upquote.sty}{\usepackage{upquote}}{}
\IfFileExists{microtype.sty}{% use microtype if available
  \usepackage[]{microtype}
  \UseMicrotypeSet[protrusion]{basicmath} % disable protrusion for tt fonts
}{}
\makeatletter
\@ifundefined{KOMAClassName}{% if non-KOMA class
  \IfFileExists{parskip.sty}{%
    \usepackage{parskip}
  }{% else
    \setlength{\parindent}{0pt}
    \setlength{\parskip}{6pt plus 2pt minus 1pt}}
}{% if KOMA class
  \KOMAoptions{parskip=half}}
\makeatother
\usepackage{longtable,booktabs,array}
\usepackage{caption}
\captionsetup[table]{skip=6pt}
\newcounter{none} % for unnumbered tables
\usepackage{calc} % for calculating minipage widths
% Correct order of tables after \paragraph or \subparagraph
\usepackage{etoolbox}
\makeatletter
\patchcmd\longtable{\par}{\if@noskipsec\mbox{}\fi\par}{}{}
\makeatother
\usepackage{graphicx}
\makeatletter
\newsavebox\pandoc@box
\newcommand*\pandocbounded[1]{% scales image to fit in text height/width
  \sbox\pandoc@box{#1}%
  \Gscale@div\@tempa{\textheight}{\dimexpr\ht\pandoc@box+\dp\pandoc@box\relax}%
  \Gscale@div\@tempb{\linewidth}{\wd\pandoc@box}%
  \ifdim\@tempb\p@<\@tempa\p@\let\@tempa\@tempb\fi% select the smaller of both
  \ifdim\@tempa\p@<\p@\scalebox{\@tempa}{\usebox\pandoc@box}%
  \else\usebox{\pandoc@box}%
  \fi%
}
% Set default figure placement to htbp
\def\fps@figure{htbp}
\makeatother
\setlength{\emergencystretch}{3em} % prevent overfull lines
\providecommand{\tightlist}{%
  \setlength{\itemsep}{0pt}\setlength{\parskip}{0pt}}
% ============================================================
% Preambulo moderno para Tarea Práctica 1
% Basado en la plantilla original del autor.
% Se conservan las interfaces: \materia, \TituloTarea, \Tony y \titulo.
% ============================================================

% Asegura que los .sty locales se encuentren dentro de la carpeta plantilla/
\makeatletter
%\def\input@path{{plantilla/}}
\makeatother

\usepackage{fontspec}      % Si no lo tienes, añádelo antes de babel
\usepackage[spanish,es-tabla,activeacute]{babel}
\selectlanguage{spanish}
\spanishdecimal{.}

\usepackage{varios}
\usepackage{needspace}  % evita que tablas cortas se partan entre páginas
\graphicspath{{plantilla/}}
% ------------------------------------------------------------------
% Paleta Línea 12 (Línea Dorada) — sobrescribe los colores UNAM.
% Todo lo que ya usa estos nombres (portada, secciones, cajas,
% hipervínculos) hereda el dorado automáticamente.
% ------------------------------------------------------------------
\definecolor{UNAMBlue}{HTML}{006199}        % dorado principal
\definecolor{UNAMAccent}{HTML}{8A6A16}      % dorado oscuro (contraste)
\definecolor{UNAMLightBlue}{HTML}{F3E7C6}   % fondo claro dorado
\definecolor{UNAMLightAccent}{HTML}{EAD6A0} % fondo claro dorado oscuro
\definecolor{UNAMMuted}{HTML}{5B5B5B}       % gris neutro, sin cambio
\definecolor{UNAMText}{HTML}{262626}        % texto principal
\definecolor{UNAMLightGray}{HTML}{FBF9F4}   % blanco hueso

% Matemáticas modernas para XeLaTeX
\usepackage{unicode-math}
\setmathfont{TeX Gyre Termes Math}

\renewcommand{\labelitemi}{$\bullet$}
\renewcommand{\labelitemii}{$\circ$}

% ============================================================
% Identidad del documento: conservar estos comandos
% ============================================================
\newcommand{\materia}{Modelos de Supervivencia y Series de Tiempo}
\newcommand{\TituloTarea}{Tarea Practica 5}
\newcommand{\Tony}{\textsc{Anthonie Byron Ornelas Clement}}

% ============================================================
% Portada
% ============================================================



\newcommand{\titulo}{%
  \begin{minipage}{0.93\textwidth}
    \vspace{0.6cm}
    \begin{center}
      \IfFileExists{escudo.png}{%
        \includegraphics[height=7cm]{escudo.png}\par\vspace{0.55cm}
      }{%
        \IfFileExists{escudo0.png}{%
          \includegraphics[height=7cm]{escudo0.png}\par\vspace{0.55cm}
        }{}
      }
      {\sffamily\small\bfseries FACULTAD DE CIENCIAS\par}
      \vspace{0.12cm}
      {\sffamily\footnotesize\color{UNAMMuted}\materia\par}
      \vspace{0.75cm}
        \colorbox{UNAMBlue}{\parbox{0.82\textwidth}{%
        \centering
        \vspace{0.22cm}
        {\sffamily\color{white}\bfseries\fontsize{24}{29}\selectfont \TituloTarea\par}
        \vspace{0.22cm}
      }}
      \par\vspace{0.3cm}
      \textcolor{UNAMAccent}{\rule{4.5cm}{1pt}}\par
      \vfill
      %{\sffamily\small\bfseries \Tony\par}
      \vspace{0.3cm}
      {\large \Tony \par}
      \date{\today}
    \end{center}
  \end{minipage}%
}

% ============================================================
% Ajustes visuales generales
% ============================================================
\setlength{\parindent}{0pt}
\setlength{\parskip}{0.55em}

% Hyperref lo administra Pandoc/R Markdown. Solo configuramos al inicio.
\makeatletter
\AtBeginDocument{%
  \@ifpackageloaded{hyperref}{%
    \hypersetup{
      colorlinks=true,
      linkcolor=UNAMBlue,
      urlcolor=UNAMBlue,
      citecolor=UNAMBlue
    }
  }{}
}
\makeatother

% Títulos
\usepackage{titlesec}
\titleformat{\section}
  {\sffamily\Large\bfseries\color{UNAMBlue}}
  {\thesection}{0.65em}{}
  [\vspace{0.18em}\color{UNAMAccent}\titlerule]

\titleformat{\subsection}
  {\sffamily\large\bfseries\color{UNAMText}}
  {\thesubsection}{0.6em}{}

\titleformat{\subsubsection}
  {\sffamily\normalsize\bfseries\color{UNAMBlue}}
  {\thesubsubsection}{0.55em}{}

% Figuras
\captionsetup{
  font=small,
  labelfont={bf,color=UNAMBlue},
  textfont={color=UNAMMuted},
  justification=centering,
  skip=7pt
}

% ============================================================
% Entornos de comunicación para el Rmd
% ============================================================
\newtcolorbox{queobservar}{%
  enhanced,
  breakable,
  colback=UNAMLightBlue,
  colframe=UNAMBlue,
  boxrule=0pt,
  leftrule=3pt,
  arc=2pt,
  left=8pt,right=8pt,top=7pt,bottom=7pt,
  before skip=8pt,after skip=10pt,
  fonttitle=\sffamily\bfseries,
  title=¿Qué debemos observar?%
}

\newtcolorbox{notatecnica}{%
  enhanced,
  breakable,
  colback=UNAMLightGray,
  colframe=UNAMMuted,
  boxrule=0pt,
  leftrule=3pt,
  arc=2pt,
  left=8pt,right=8pt,top=7pt,bottom=7pt,
  before skip=8pt,after skip=10pt,
  fonttitle=\sffamily\bfseries,
  title=Nota técnica%
}

\newtcolorbox{hallazgo}{%
  enhanced,
  breakable,
  colback=UNAMLightAccent,
  colframe=UNAMAccent,
  boxrule=0pt,
  leftrule=4pt,
  arc=2pt,
  left=8pt,right=8pt,top=8pt,bottom=8pt,
  before skip=9pt,after skip=11pt,
  fonttitle=\sffamily\bfseries,
  title=Hallazgo principal%
}


\newtcolorbox{citadestacada}{%
  enhanced,
  breakable,
  colback=white,
  colframe=UNAMBlue,
  boxrule=0pt,
  leftrule=2.4pt,
  arc=0pt,
  left=12pt,right=10pt,top=6pt,bottom=6pt,
  before skip=8pt,after skip=10pt,
  fontupper=\itshape\color{UNAMText}
}
\usepackage{booktabs}
\usepackage{longtable}
\usepackage{array}
\usepackage{multirow}
\usepackage{wrapfig}
\usepackage{float}
\usepackage{colortbl}
\usepackage{pdflscape}
\usepackage{tabularx}
\usepackage{xltabular}
\usepackage{threeparttable}
\usepackage{threeparttablex}
\usepackage[normalem]{ulem}
\usepackage{makecell}
\usepackage{xcolor}
\usepackage{bookmark}
\IfFileExists{xurl.sty}{\usepackage{xurl}}{} % add URL line breaks if available
\urlstyle{same}
% fallback for those not using the hyperref driver hyperxmp:
\makeatletter
\@ifundefined{xmpquote}{\newcommand{\xmpquote}[1]{#1}}{}
\makeatother
\hypersetup{
  hidelinks,
  pdfcreator={LaTeX via pandoc}}

\title{\titulo}
\author{}
\date{\vspace{-2.5em}}

\begin{document}
\maketitle

\section{Introducción}\label{introducciuxf3n}

La empresa que nos contrató opera una tienda en línea y quiere saber qué
modificación a su producto mejoraría la rentabilidad. El análisis de la
información recolectada en el último año (1,000 clientes)
arroja un hallazgo incómodo pero accionable: \textbf{ningún rasgo observable del
cliente anticipa su fuga} --- edad, ingreso, gasto, devoluciones, soporte,
satisfacción y respuesta a promociones discriminan igual entre quienes abandonan
y quienes no. El riesgo es de \textbf{52.6\% en el año, presente
todo el año con aceleración al final} (según la curva de Kaplan--Meier, la
mitad de la base habría fugado hacia el día
269).

La consecuencia directa para el negocio: esperar señales para retener no
funciona con estos datos, así que la acción recomendada es un \textbf{ciclo de
re-engagement universal y calendario fijo} (contacto preventivo antes del día
180, cuando la permanencia aún es del
70.4\%), más una \textbf{re-instrumentación del registro de clientes}
para que el siguiente análisis sí pueda detectar segmentos de riesgo.

\section{Planteamiento del problema y propuesta}\label{planteamiento-del-problema-y-propuesta}

La tienda vende en línea a clientes cuya actividad se mide por compras; el
producto incluye promociones, atención a soporte y comunicación por correo. La
dirección pregunta qué cambiar en el producto para retener más clientes: cada
punto de retención preserva el ingreso que el cliente genera sin costo de
adquisición.

En lugar de clasificar clientes como ``fugados o no'' al final del año, usamos
\textbf{análisis de supervivencia}, que estudia \emph{cuánto tarda en ocurrir un evento}.
Aquí el evento es la fuga (el cliente marcado como abandonado) y el tiempo es
la inactividad: los días transcurridos desde su última compra. Esta técnica
aporta dos ventajas directas para decidir: entrega una \textbf{línea de tiempo del
riesgo} (en qué días conviene actuar) y \textbf{compara grupos} (¿retiene más quien
responde promociones?). Un cliente que no ha fugado al cierre del análisis
simplemente aporta la información de que sobrevivió hasta su último registro
(censura); no se le descarta.

Las preguntas que organiza el análisis:

\begin{enumerate}
\def\labelenumi{\arabic{enumi}.}
\tightlist
\item
  ¿En qué momento del año se concentra la fuga? (timing de la acción.)
\item
  ¿Algún segmento de clientes fuga más rápido que otro? (¿a quién apuntar?)
\item
  ¿Las promociones y el soporte están asociados con la retención? (¿qué
  modificar del producto?)
\end{enumerate}

\section{La base y su limpieza}\label{la-base-y-su-limpieza}

La base tiene 1,000 clientes y 15 variables (identificador,
perfil demográfico y económico, comportamiento de compra y abandono). La
inspección y las decisiones se resumen en la Tabla \ref{tab:limpieza}; \textbf{no se
eliminó ningún cliente} y ninguna variable requirió imputación o truncado: no
hay duplicados, no hay perdidos y ninguna de las 10 variables numéricas tiene
observaciones fuera de 1.5 rangos intercuartílicos (suma de extremos detectados:
0; Figura \ref{fig:extremos}).

\begin{table}[!h]
\centering
\caption{\label{tab:limpieza}Decisiones de limpieza y construcción de variables (antes → después).}
\centering
\fontsize{9}{11}\selectfont
\begin{tabular}[t]{m{1.8cm}m{2.2cm}p{9.4cm}|m{1.8cm}m{2.2cm}p{9.4cm}|m{1.8cm}m{2.2cm}p{9.4cm}}
\hline
\cellcolor[HTML]{264653}{\textcolor{white}{\textbf{Paso}}} & \cellcolor[HTML]{264653}{\textcolor{white}{\textbf{Decisión}}} & \cellcolor[HTML]{264653}{\textcolor{white}{\textbf{Detalle}}}\\
\hline
\cellcolor[HTML]{F1FAEE}{Lectura} & \cellcolor[HTML]{F1FAEE}{CSV crudo → datos} & \cellcolor[HTML]{F1FAEE}{15 columnas tal como se reciben}\\
\hline
\cellcolor[HTML]{F1FAEE}{Duplicados} & \cellcolor[HTML]{F1FAEE}{Verificación} & \cellcolor[HTML]{F1FAEE}{0 duplicados → sin acción}\\
\hline
\cellcolor[HTML]{F1FAEE}{Perdidos} & \cellcolor[HTML]{F1FAEE}{Verificación} & \cellcolor[HTML]{F1FAEE}{0 perdidos → sin imputación}\\
\hline
\cellcolor[HTML]{F1FAEE}{Extremos} & \cellcolor[HTML]{F1FAEE}{Verificación} & \cellcolor[HTML]{F1FAEE}{0 fuera de 1.5×IQR en las 10 numéricas → sin truncado}\\
\hline
\cellcolor[HTML]{F1FAEE}{Categóricas} & \cellcolor[HTML]{F1FAEE}{Recodificación} & \cellcolor[HTML]{F1FAEE}{Gender y Promotion\_Response a factor; Email\_Opt\_In y Target\_Churn de texto «True/False» a lógico}\\
\hline
\cellcolor[HTML]{F1FAEE}{Evento} & \cellcolor[HTML]{F1FAEE}{Construcción} & \cellcolor[HTML]{F1FAEE}{Target\_Churn → evento (1 fuga, 0 retención): 526 fugas (52.6\%)}\\
\hline
\cellcolor[HTML]{F1FAEE}{Tiempo} & \cellcolor[HTML]{F1FAEE}{Construcción} & \cellcolor[HTML]{F1FAEE}{Last\_Purchase\_Days\_Ago → días de inactividad (1–364)}\\
\hline
\cellcolor[HTML]{F1FAEE}{Satisfacción} & \cellcolor[HTML]{F1FAEE}{Agrupación} & \cellcolor[HTML]{F1FAEE}{Satisfaction\_Score 1–2 → «Baja» (394); 3–5 → «Alta» (606)}\\
\hline
\end{tabular}
\end{table}

La construcción del \textbf{par de supervivencia} merece nota explícita: el \emph{evento}
es \texttt{Target\_Churn} (1 = fugado) y el \emph{tiempo} es \texttt{Last\_Purchase\_Days\_Ago}
(inactividad al cierre del año). Los clientes sin fuga quedan \textbf{censurados}:
sabemos que seguían activos hasta su última compra, entre los días
1 y 364 del año, pero no su vida completa.
La limitación que esto impone --- la fuga se mide contra la última compra y no
contra una fecha de baja declarada --- se retoma en las conclusiones.

\section{¿Quiénes son y cómo compran? (exploración)}\label{quiuxe9nes-son-y-cuxf3mo-compran-exploraciuxf3n}

El perfil del cliente es amplio y balanceado: edades de
18 a 69 años, ingresos anuales de
20 a 200 miles de dólares,
antigüedad de 1 a 19 años y
género repartido en tres tercios (Figura \ref{fig:perfil}). El comportamiento
de compra es igualmente disperso y sin valores extremos: de
1 a 99 compras al año,
gasto total de 109 a 10,000 dólares
y transacciones promedio de 10 a
500 dólares (Figura \ref{fig:extremos}).

\begin{figure}

{\centering \includegraphics{Tarea_Practica_5_files/figure-latex/perfil-1} 

}

\caption{Perfil de los clientes: ninguna característica demográfica concentra a la base; el público es heterogéneo y balanceado.}\label{fig:perfil}
\end{figure}

\begin{figure}

{\centering \includegraphics{Tarea_Practica_5_files/figure-latex/extremos-1} 

}

\caption{Diagramas de caja del comportamiento de compra: no hay valores extremos fuera de 1.5×IQR en ninguna variable; no se requirió truncado ni winsorización.}\label{fig:extremos}
\end{figure}

Las variables numéricas entre sí prácticamente no se correlacionan
(Figura \ref{fig:correlaciones}): no existe un ``cliente grande'' (mucho ingreso
y mucho gasto, muchas compras) ni combinaciones de variables que agrupen a la
base. Esta primera señal --- la base no tiene segmentos naturales --- anticipa el
hallazgo central del análisis de supervivencia.

\begin{figure}

{\centering \includegraphics{Tarea_Practica_5_files/figure-latex/correlaciones-1} 

}

\caption{Correlaciones entre variables numéricas: todas cercanas a cero; la base no tiene segmentos naturales.}\label{fig:correlaciones}
\end{figure}

\section{¿Qué distingue a quien fuga? (comparación inicial)}\label{quuxe9-distingue-a-quien-fuga-comparaciuxf3n-inicial}

Al cruzar cada variable con el abandono (Figura \ref{fig:churn-rel}) no
aparece ningún contraste: la distribución de satisfacción, contactos a soporte
y la respuesta a promociones es la misma entre fugados y activos. La fracción de fuga oscila
\textbf{sin tendencia entre 48.3\% y
60.8\%} en todas las ventanas del año.

\begin{figure}

{\centering \includegraphics{Tarea_Practica_5_files/figure-latex/churn-rel-1} 

}

\caption{Fuga frente a satisfacción, soporte y promociones: las distribuciones de fugados (rojo) y activos (verde) son indistinguibles en toda variable.}\label{fig:churn-rel}
\end{figure}

\section{¿Cuándo se van los clientes? (análisis de supervivencia)}\label{cuuxe1ndo-se-van-los-clientes-anuxe1lisis-de-supervivencia}

La curva de Kaplan--Meier global (Figura \ref{fig:km-global}) responde la
primera pregunta: la fuga \textbf{no se concentra en ningún periodo}. El abandono se
acumula a un ritmo promedio de 1.45 clientes por día: el
riesgo es plano durante la mayor parte del año y se acelera al final, en la
cola de la distribución de inactividad (Figura
\ref{fig:km-global}). La permanencia es del
94.9\% al día 30, 85.9\% al día 90 y
70.4\% al día 180; la mediana de ``sobrevivencia'' es de
269 días (IC 95\%: 255--283),
es decir, la mitad de la base habría fugado hacia los 9 meses de inactividad.
Parte de la aceleración final de la curva refleja el cierre de la
ventana de observación (un cliente solo puede acumular hasta 364 días de
inactividad) y no un cambio de comportamiento.

\begin{figure}

{\centering \includegraphics{Tarea_Practica_5_files/figure-latex/km-global-1} 

}

\caption{Kaplan–Meier global: la permanencia cae de manera pareja la mayor parte del año y se acelera al final (cola de inactividad); la mitad de la base habría fugado hacia el día 269. La acción de retención debe activarse antes del día 180.}\label{fig:km-global}
\end{figure}

Las comparaciones por segmento (Figura \ref{fig:km-segmentos}) confirman la
segunda pregunta. Para la satisfacción se prueba H0: la permanencia es igual
entre clientes con satisfacción baja y alta, contra la alternativa de que
difiere; las curvas de satisfacción baja (394 clientes) y alta
(606) se traslapan (log-rank p = 0.230). La misma
prueba entre quienes respondieron, ignoraron o cancelaron su suscripción a las
promociones tampoco detecta diferencias (p = 0.509). Ninguna
diferencia es estadísticamente detectable.

\begin{figure}

{\centering \includegraphics{Tarea_Practica_5_files/figure-latex/km-segmentos-1} 

}

\caption{Permanencia según satisfacción y respuesta a promociones: las curvas se traslapan (log-rank p=0.23 y p=0.51); ningún segmento retiene mejor que otro.}\label{fig:km-segmentos}
\end{figure}

Para ponderar todas las variables a la vez, ajustamos el modelo de Cox, que
estima cuánto multiplica cada característica el riesgo de fuga manteniendo el
resto constante (Tabla \ref{tab:cox-tabla}). La lectura ejecutiva de la tabla:
\textbf{ningún hazard ratio se aleja de 1.0 con confianza} --- el intervalo de cada
variable incluye al 1.0 --- y la significancia más cercana (gasto total,
p = 0.054) no alcanza el umbral convencional de 0.05. Para
validar las estimaciones se prueba H0: los riesgos proporcionales se mantienen
en el tiempo (supuesto del modelo); \textbf{no se detecta violación del supuesto}
(prueba global p = 0.77), de modo que las estimaciones son
válidas; lo que dicen es simplemente que no hay efectos.

\begin{table}[!h]
\centering
\caption{\label{tab:cox-tabla}Modelo de Cox: hazard ratios (HR) por variable. HR > 1 aumenta el riesgo de fuga; todo intervalo contiene al 1.0 → ningún efecto detectable.}
\centering
\fontsize{9}{11}\selectfont
\begin{tabular}[t]{c|c|c|c}
\hline
\cellcolor[HTML]{264653}{\textcolor{white}{\textbf{Variable}}} & \cellcolor[HTML]{264653}{\textcolor{white}{\textbf{HR}}} & \cellcolor[HTML]{264653}{\textcolor{white}{\textbf{IC}}} & \cellcolor[HTML]{264653}{\textcolor{white}{\textbf{p}}}\\
\hline
\cellcolor[HTML]{F1FAEE}{Edad} & \cellcolor[HTML]{F1FAEE}{1.000} & \cellcolor[HTML]{F1FAEE}{[0.994, 1.005]} & \cellcolor[HTML]{F1FAEE}{0.933}\\
\hline
\cellcolor[HTML]{F1FAEE}{Ingreso anual (miles USD)} & \cellcolor[HTML]{F1FAEE}{0.998} & \cellcolor[HTML]{F1FAEE}{[0.997, 1.000]} & \cellcolor[HTML]{F1FAEE}{0.064}\\
\hline
\cellcolor[HTML]{F1FAEE}{Gasto por cada 1,000 USD} & \cellcolor[HTML]{F1FAEE}{1.031} & \cellcolor[HTML]{F1FAEE}{[1.000, 1.064]} & \cellcolor[HTML]{F1FAEE}{0.054}\\
\hline
\cellcolor[HTML]{F1FAEE}{Antigüedad (años)} & \cellcolor[HTML]{F1FAEE}{0.992} & \cellcolor[HTML]{F1FAEE}{[0.977, 1.008]} & \cellcolor[HTML]{F1FAEE}{0.323}\\
\hline
\cellcolor[HTML]{F1FAEE}{Compras al año} & \cellcolor[HTML]{F1FAEE}{1.000} & \cellcolor[HTML]{F1FAEE}{[0.997, 1.003]} & \cellcolor[HTML]{F1FAEE}{0.795}\\
\hline
\cellcolor[HTML]{F1FAEE}{Transacción por cada 1,000 USD} & \cellcolor[HTML]{F1FAEE}{0.709} & \cellcolor[HTML]{F1FAEE}{[0.394, 1.276]} & \cellcolor[HTML]{F1FAEE}{0.251}\\
\hline
\cellcolor[HTML]{F1FAEE}{Devoluciones} & \cellcolor[HTML]{F1FAEE}{0.987} & \cellcolor[HTML]{F1FAEE}{[0.958, 1.017]} & \cellcolor[HTML]{F1FAEE}{0.389}\\
\hline
\cellcolor[HTML]{F1FAEE}{Contactos a soporte} & \cellcolor[HTML]{F1FAEE}{1.031} & \cellcolor[HTML]{F1FAEE}{[0.970, 1.096]} & \cellcolor[HTML]{F1FAEE}{0.332}\\
\hline
\cellcolor[HTML]{F1FAEE}{Satisfacción (1–5)} & \cellcolor[HTML]{F1FAEE}{1.000} & \cellcolor[HTML]{F1FAEE}{[0.940, 1.064]} & \cellcolor[HTML]{F1FAEE}{0.997}\\
\hline
\cellcolor[HTML]{F1FAEE}{Aceptó correos (sí vs no)} & \cellcolor[HTML]{F1FAEE}{0.932} & \cellcolor[HTML]{F1FAEE}{[0.784, 1.108]} & \cellcolor[HTML]{F1FAEE}{0.425}\\
\hline
\cellcolor[HTML]{F1FAEE}{Respondió promo (vs ignoró)} & \cellcolor[HTML]{F1FAEE}{1.045} & \cellcolor[HTML]{F1FAEE}{[0.842, 1.297]} & \cellcolor[HTML]{F1FAEE}{0.691}\\
\hline
\cellcolor[HTML]{F1FAEE}{Canceló promo (vs ignoró)} & \cellcolor[HTML]{F1FAEE}{0.938} & \cellcolor[HTML]{F1FAEE}{[0.757, 1.162]} & \cellcolor[HTML]{F1FAEE}{0.557}\\
\hline
\end{tabular}
\end{table}

\section{Qué implica esto para el negocio}\label{quuxe9-implica-esto-para-el-negocio}

\textbf{1. El riesgo es plano y alto: retener requiere actuar a todos, no a unos
pocos.} Con 52.6\% de fuga anual, sin segmentos
distinguibles y con un riesgo que solo se acelera al final del año, un programa de retención dirigido a
``clientes en riesgo'' no tiene a dónde apuntar: las señales disponibles no
anticipan nada. La modificación al producto que este análisis sostiene es un
\textbf{ciclo de re-engagement universal}: contacto preventivo a toda la base antes
del día 180 (cuando la permanencia sigue en 70.4\%) --- por
ejemplo, incentivo de recompra y recordatorio de valor --- repetido
trimestralmente.

\textbf{2. Las promociones actuales no muestran efecto medible; no conviene
escalarlas.} Fugas y permanencia se comportan igual entre quienes respondieron
y quienes cancelaron su suscripción a la campaña
(p = 0.509; HR de 1.04, IC 95\%
\texttt{{[}0.84,\ 1.30{]}}). Antes de
invertir más en la campaña, hay que rediseñarla para que sea capaz de mover la
retención.

\textbf{3. Re-instrumentar el registro de clientes es la modificación de mayor valor
a largo plazo.} Este análisis se limita por cómo se capturó la información: no
hay fecha de alta, ni historial de compras fechado, ni fecha de baja --- solo un
corte de ``días desde la última compra'' y una marca final de abandono. Registrar
esas fechas convertiría el siguiente análisis en un seguimiento real del
tiempo de vida del cliente, donde sí puedan detectarse segmentos de riesgo y
medirse las campañas de retención.

\textbf{Nota técnica (para quien domina el tema).} Tiempo = \texttt{Last\_Purchase\_Days\_Ago}
(1--364 días), evento = \texttt{Target\_Churn}
(526 fugas en 1,000 clientes); censurados a la derecha
sin fuga. Cox con todas las covariables de la Tabla 2: sin HR significativo;
\texttt{cox.zph} global p = 0.77 (sin violación del supuesto de
riesgos proporcionales). La base muestra patrones compatibles con datos simulados
(varianzas planas, nulas correlaciones), por lo que las conclusiones de negocio
son provisionales mientras la instrumentación no mejore.

\end{document}
